% !TEX root = ../notes.tex

\documentclass[../notes.tex]{subfiles}

\begin{document}

\section{March 9}
Here we go.

\subsection{The Bruhat--Tits Tree}
Let's recall our story in some level of generality. Let $G$ be a split reductive group over a nonarchimedean local field $F$, and choose a maximal torus $T\subseteq G$ contained in a Borel $B\subseteq G$. For any unramified character $\chi$, there is a unique spherical subquotient of $\op{Ind}_B^G\chi$, so one accordingly finds that spherical representations are parameterized such unramified characters up to the $W$-action.

On the other hand, spherical representations are simply given by simple $\mc H_K$-modules. Thus, we would like to understand our Hecke algebra $\mc H_K=C_c(K\backslash G/K)\,dg$.
\begin{lemma} \label{lem:k-coset-to-lattice}
	Consider the group $G\coloneqq\op{GL}_n$ over a nonarchimedean local field $F$. Then $K\backslash G/K$ is in bijection with the set of $G$-orbits of pairs $\OO$-lattices $(\Lambda_1,\Lambda_2)$ in $F^n$.
\end{lemma}
\begin{proof}
	Send $KgK$ to $(\OO^n,g\OO^n)$; this is of course $K$-invariant on the right, and it is $K$-invariant on the left because $(\OO^n,g\OO^n)$ and $(\OO^n,kg\OO^n)$ are in the same $G$-orbit. Conversely, given any pair $(\Lambda_1,\Lambda_2)$, one can find a distinguished element in the $G$-orbit with $\Lambda_1=\OO^n$ so that $(\Lambda_1,\Lambda_2)=(\OO^n,g\OO^n)$ for some $g\in G$, and the element $g\in G$ is well-defined in $K\backslash G/K$.
\end{proof}
Thus, we are motivated to understand lattices.
\begin{proposition} \label{prop:classify-lattice-pairs}
	Fix two lattices $\Lambda_1,\lambda_2\subseteq F^n$. Then there is a basis $\{e_1,\ldots,e_n\}\subseteq F^n$ of $\Lambda_1$ so that
	\[\Lambda_2=\varpi^{d_1}\OO e_1\oplus\cdots\oplus\varpi^{d_n}\OO e_n\]
	for some decreasing sequence $\{d_1,\ldots,d_n\}$ of integers. In fact, the sequence of integers does not depend on the choice of basis.
\end{proposition}
\begin{proof}
	We induct on $n$. In the case where $n=0$, the statement is vacuous, and for $n=1$, the statement has no content. By rescaling the conclusion, we may assume that $\Lambda_1\subseteq\Lambda_2$ but $\Lambda_1\not\subseteq\varpi\Lambda_2$. Then the natural inclusion induces a nonzero map
	\[\Lambda_1/\varpi\Lambda_1\to\Lambda_2/\varpi\Lambda_2.\]
	As such, we just choose a basis $\{e_1,\ldots,e_r\}$ of the image and lift them to $\Lambda_1$, where they will continue to be linearly independent, and they generate a saturated sublattice $\Lambda_0$ of both $\Lambda_1$ and $\Lambda_2$. Thus, we may set $d_1=\cdots=d_r=0$, and the remaining integers are found by induction by passing to the pair of lattices $(\Lambda_1/\Lambda_0,\Lambda_2/\Lambda_0)$ of smaller rank. Indeed, one can lift the basis for $\Lambda_1/\Lambda_0$ to the required remaining basis vectors $\{e_{r+1},\ldots,e_n\}$; one will need to modify these $e_i$s by a vector in $\Lambda_0$ to ensure that $\varpi^{d_i}e_i\in\Lambda_2$ (because a priori we merely have $\varpi^{d_i}e_i\in\Lambda_0$).
\end{proof}
\begin{remark}
	This is basically equivalent to the Smith normal form: we are trying to use integral row and column operations to make a given matrix into a diagonal one.
\end{remark}
\begin{corollary}[Cartan decomposition] \label{cor:cartan-decomp}
	Consider the group $G\coloneqq\op{GL}_n$ over a nonarchimedean local field $F$. There is a disjoint union
	\[G=\bigsqcup_{d_1\ge\cdots\ge d_n}K\op{diag}\left(\varpi^{d_1},\ldots,\varpi^{d_n}\right)K.\]
\end{corollary}
\begin{proof}
	Combine \Cref{lem:k-coset-to-lattice} with \Cref{prop:classify-lattice-pairs}.
\end{proof}
\begin{remark}
	The same argument shows that
	\[\op{PGL}_2(F)=\bigsqcup_{d\ge0}K\op{diag}\left(\varpi^d,1\right)K.\]
	The corresponding statement of \Cref{prop:classify-lattice-pairs} simply needs to notice that if we replace the lattices with homothety classes, then we are looking at pairs of integers $(d_1,d_2)$ up to shifting, which can be uniquely presented as $(d,0)$. The Bruhat--Tits tree is now recovered as having vertices given by the homothety classes of lattices and edges which are connected if and only if $([\Lambda],[\Lambda'])$ is represented by the pair of integers $(1,0)$. Accordingly, general pairs $([\Lambda],[\Lambda'])$ classified by $(d,0)$ have distance $d$.
\end{remark}
\begin{remark}
	For $\op{PGL}_n$, the Bruhat--Tits tree becomes a Bruhat--Tits building is a simplicial complex of dimension $(n-1)$, where vertices are given by lattices $\Lambda$.
\end{remark}
Thus, we see that $\mc H_K$ admits a basis of indicators of the double cosets of \Cref{cor:cartan-decomp}.

\subsection{The Spherical Hecke Algebra}
Here is the usual application of the Cartan decomposition.
\begin{lemma}[Gelfand] \label{lem:gelfand-trick}
	Consider the group $G\coloneqq\op{GL}_n$ over a nonarchimedean local field $F$. Then $\mc H_K$ is commutative.
\end{lemma}
\begin{proof}
	We use the Gelfand trick: consider the map $G\to G$ given by the transpose. This operation induces an anti-commutative map $C(G)\to C(G)$ given by sending $f$ to $f^\intercal(g)\coloneqq f(g^\intercal)$. In particular, we see that $(-)^\intercal$ is linear, but $(f_1f_2)^\intercal=f_2^\intercal f_1^\intercal$. However, $(-)^\intercal$ fixes the basis elements
	\[1_{K\op{diag}\left(\varpi^{d_1},\ldots,\varpi^{d_n}\right)K}\]
	of \Cref{cor:cartan-decomp} are fixed by $(-)^\intercal$, so $(-)^\intercal$ is the identity. We conclude that $\mc H_K$ is commutative.
\end{proof}
The above miraculous proof turns out to work for general split reductive groups.
\begin{proposition}[Cartan decomposition]
	Fix a split reductive group $G$ over a nonarchimedean local field $F$, and let $T\subseteq G$ be a split maximal torus. For $\lambda\in X_*(T)$, we let $\varpi^\lambda$ denote $\lambda(\varpi)\in T(F)$. Then
	\[G(F)=\bigsqcup_{\text{dominant }\lambda\in X_*(T)}G(\OO)\varpi^\lambda G(\OO).\]
	Here, a coweight $\lambda$ is dominant if and only if $\langle\lambda,\alpha_i\rangle\ge0$ for all simple roots $\alpha_i$.
\end{proposition}
\begin{proof}
	Omitted!
\end{proof}
\begin{theorem}
	Fix a split reductive group $G$ over a nonarchimedean local field $F$. Then the Hecke algebra $\mc H_{G(\OO)}$ is commutative.
\end{theorem}
\begin{proof}
	It is a fact that there is an algebraic group involution $\sigma$ of $G$ for which $\sigma(t)=t^{-1}$ for all $t\in T$; it also sends a choice of Borel to the opposite Borel. (In fact, this $\sigma$ is defined over $\ZZ$. Notably, the action by the longest Weyl element should work.) Then one replaces $(-)^\intercal$ in the proof of \Cref{lem:gelfand-trick} with the function $\tau\colon G(F)\to G(F)$ given by $g\mapsto\sigma(g)^{-1}$, and the entire argument goes through.
\end{proof}
Of course, we are still interested in the ring structure of the spherical Hecke algebra. For example, with $G=\op{GL}_n$, we claim that
\[1_{K\op{diag}\left(\varpi^{d_1},\ldots,\varpi^{d_n}\right)K}*1_{K\op{diag}\left(\varpi^{d_1'},\ldots,\varpi^{d_n'}\right)K}=1_{K\op{diag}\left(\varpi^{d_1+d_1'},\ldots,\varpi^{d_n+d_n'}\right)K}+\text{lower-order terms}.\]
(Here, our collection of decreasing sequences of integers is given by the usual order on partitions.\footnote{One can show that having $d\ge d'$ is equivalent to $K\varpi^d K$ containing $K\varpi^{d'}$ in its closure.}) Indeed, it is enough to compute the left-hand side on matrices of the form $\op{diag}(\varpi^{\ell_1},\ldots,\varpi^{\ell_n})$. By expanding out the convolution, we see that this value is given by the number of lattices $\Lambda$ which can fit in a sequence
\[\OO^{\oplus n}\stackrel d\to\Lambda\stackrel{d'}\to\varpi^{\ell_1}\OO\oplus\cdots\oplus\varpi^{\ell_n}\OO,\]
where $\stackrel d\to$ means that the pair of consecutive lattices is parameterized by $d$.
\begin{example}
	We go down to $\op{PGL}_2$. Then we are trying to show
	\[1_{K\op{diag}\left(\varpi^d,1\right)K}*1_{K\op{diag}\left(\varpi^{d'},1\right)K}=1_{K\op{diag}\left(\varpi^{d+d'},1\right)K}+\text{lower-order terms}.\]
	Well, on $\op{diag}\left(\varpi^\ell,1\right)$, we are asking for the number of lattices $\Lambda$ which are a distance of $d$ from $\OO^{\oplus2}$ and a distance $d'$ from $\varpi^\ell\OO\oplus\OO$. For example, if $d+d'<\ell$, there are no such lattices; if $d+d'=\ell$, there is in fact a distinguished $\Lambda$ for $\ell=d+d'$. However, there are allowed to be ``lower-order terms'' for $\ell>d+d'$.
\end{example}
Admitting the claim, we conclude that $\op{gr}\mc H_K$ is a polynomial ring on $X_*(T)/W$. This sort of discussion generalizes to prove the following.
\begin{theorem}
	Fix a split reductive group $G$ over a nonarchimedean local field $F$, and set $K\coloneqq G(\OO)$. Then $\op{gr}\mc H_K$ is isomorphic to $\CC[X_*(T)^{\mathrm{dom}}]$, where $X_*(T)^{\mathrm{dom}}$ denotes the monoid of dominant weights.
\end{theorem}
\begin{proof}
	Proceed as above with more effort.
\end{proof}
\begin{remark}
	It follows purely formally that $\mc H_K$ (with some Five lemma argument) that $\mc H_K$ is a polynomial ring of rank $\op{rank}G$.
\end{remark}

\subsection{The Satake Isomorphism, Again}
It may be desirable to have a more explicit description of $\mc H_K$, without the need to pass to the associated graded. This is given by the Satake isomorphism.
\begin{theorem}[Satake]
	Fix a split reductive group $G$ over a nonarchimedean local field $F$, and set $K\coloneqq G(\OO)$. Then there is a canonical isomorphism
	\[\mc H_K=\CC[X_*(T)]^W.\]
\end{theorem}
\begin{example}
	For $G=\op{PGL}_2$, one sees that $\mc H_K$ is a polynomial ring with generator $t\coloneqq1_{K\op{diag}(\varpi,1)K}$. On the other hand, $\CC[X_*(T)]^W$ is just $\CC[\ZZ x]^{S_2}$, which is $\CC\left[x+x^{-1}\right]$. It turns out that the Satake isomorphism sends $t$ to a scalar times $x+x^{-1}$.
\end{example}
\begin{proof}[Sketch]
	Let's start by defining the map $S\colon\mc H_K\to\CC[X_*(T)]^W$. Thus, given $\varphi\in\mc H_K$, we are being asked for a coefficient $S\varphi(\lambda)$ for each $\lambda\in X_*(T)$. To this end, we take as definition
	\[S\varphi(\lambda)\coloneqq\delta_B^{1/2}(\varpi^\lambda)\int_{N(F)}\varphi\left(\varpi^\lambda x\right)\,dx,\]
	where $B\subseteq G$ is a chosen Borel containing $T$, and $N\subseteq B$ is the unipotent radical. One can check that $S\varphi(w\lambda)=S\varphi(\lambda)$ always; this is not hard to see directly for $G=\op{GL}_2$, and then the case for general $\op{GL}_n$ can probably be checked via transpositions.

	It is rather non-obvious that this map has any good properties, so let's give another construction of $S$. Recall that $\op{Ind}_B^G\chi$ basically consists of $\chi$-eigenvectors for functions $N(F)\backslash G(F)\to\CC$. Because $\chi$ is unramified, we thus see that our functions in fact descend to $T(\OO)N(F)\backslash G(F)\to\CC$; note that the quotient continues to admit a natural action by $X_*(T)=T(F)/T(\OO)$. Thus, as $\chi$ varies, we find the ``universal'' unramified principal series
	\[\op{Ind}^{G(F)}_{T(\OO)N(F)}1.\]
	If we take $K$-invariants, then we retain the $X_*(T)$-action and exchange the $G$-action for an $\mathcal H_K$-action, and the space is given by those functions $T(\OO)N(F)\backslash G(F)/K\to\CC$, but the Iwasawa decomposition tells us that this latter set is given by $X_*(T)$. We thus have a bimodule
	\[M_{\mathrm{sph}}\coloneqq C_c(T(\OO)N(F)\backslash G(F)/K)\]
	with commuting actions by $\mc H_K$ and $\CC[X_*(T)]$. Because the two actions commute, we receive a map $\mc H_K\to\op{End}_{\CC[X_*(T)]}(M_{\mathrm{sph}})$, but $M_{\mathrm{sph}}$ is free of rank $1$ over $\CC[X_*(T)]$, so we are actually receiving a map
	\[\mc H_K\to\op{End}_{\CC[X_*(T)]}M_{\mathrm{sph}}=\CC[X_*(T)].\]
	We can also use this description to explain the $W$-invariance. For generic $\chi$, one has that $\op{Ind}_B^G\chi=\op{Ind}_B^Gw\chi$ by some intertwining operators, which descends to a $W$-action on $C_c(T(\OO)N(F)\backslash G(F)/K)$.

	It remains to show that we actually have an isomorphism. Well, it is enough to show that we have an isomorphism of vector spaces, for which one can pass to some associated graded things, so the calculation is similar to what was done before to show that $\mc H_K$ is isomorphic to a polynomial ring.
\end{proof}
Let's put this all together. We see that spherical representations are in bijection with simple $\mc H_K$-modules, which are in bijection with maximal ideals of $\mc H_K$, which are in bijection with maximal ideals of $\CC[X_*(T)]^W$, which are in bijection with $W$-orbits in $\widehat T=\op{Hom}(X_*(T),\CC^\times)$. Thus, our spherical representations are in bijection with $\widehat T/W$.
\begin{remark}
	The group $\widehat T$ is known as its ``Langlands dual.'' Note that the functor $T\mapsto\widehat T$ (from the category of tori) is contravariant, which is one explanation for why this is a dual.
\end{remark}

\end{document}